\SourceLeaf{146}
neuhinzukommenden Ehepaare, als das Verhältniss, nach welchem in einer solchen Gruppe die verschiedenen Altersstufen gemischt sind, Jahr für Jahr sich gleich bleibt, imgleichen, dass das Absterben genau mit den Mortalitätstafeln gleichen Schritt hält, so wird jener Beharrungszustand seinen Namen nach aller Strenge verdienen; sowohl die Anzahl der Ehepaare in der Gesellschaft wird ganz ungeändert bleiben, als die Anzahl der Witwen, nachdem diese ihren höchsten Werth einmal erreicht hat. Es ist ferner klar, dass, wenn man sich eine zweite ähnliche Gesellschaft vorstellt, in welcher die neubeitretenden genau in demselben Verhältnisse gemischt sind, wie in der ersten, welche aber doppelt so viele Ehepaare umfasst, die höchste Witwenzahl auch doppelt so gross sein werde; oder, um es allgemein auszudrücken, das Verhältniss der stehenden Ehen zu der Zahl der Witwen, nach eingetretenem Beharrungszustande, wird nicht von dem Umfange der Association (insofern er nur als constant bleibend betrachtet wird), sondern bloss von dem Verhältnisse abhängen, nach welchem die verschiedenen Altersstufen der neu beitretenden gemischt sind, und, wenn letzteres Verhältniss vollkommen bekannt wäre, würde ersteres sich a priori durch Rechnung bestimmen lassen. Eintreten wird übrigens der Beharrungszustand, wo nicht früher, doch jedenfalls dann, wenn die Stammtheilnehmer alle ausgestorben sind, was, wenn die extremsten Fälle berücksichtigt werden sollen, möglicherweise sich bis 80 Jahre nach dem ersten Zusammentreten verzögern könnte. Indessen kann man annehmen, dass schon nach 45—50 Jahren der wirkliche Zustand dem Beharrungszustande sehr nahe gekommen sein wird\footnote{\textsc{Kritter} begnügt sich, in seinem Gutachten, schon mit 40 Jahren.}. Um die Vorstellungen mehr zu fixiren, will ich beispielshalber bestimmte Zahlen nennen, welche jedoch, da sie nur zur Erläuterung des sonst abstracten Vortrags dienen sollen, auf vollkommen scharfe Angemessenheit keinen Anspruch machen. Die Gesellschaft bestehe aus 2600 Ehepaaren, zu denen jährlich 130 neue hinzutreten, während durchschnittlich eben so viele abgehen, und zwar 70 durch den Tod des Mannes, 60 durch den Tod der Frau. Von den während eines Jahres entstehenden 70 Witwen stirbt eine schon in demselben Jahre, so dass am Schluss des ersten Jahres 69 Witwen vorhanden sind. Von diesen sind am Ende des zweiten Jahres weitere 2 verstorben, dagegen wieder 69 neue Witwen hinzugekommen, folglich zusammen 136 vorhanden. Von diesen sterben während des dritten Jahrs 4, und indem die neu hinzugekommenen wieder 69 betragen, ist die Gesammtzahl am Schluss des dritten Jahres 201. Auf diese Weise immer langsamer fortschreitend mag die Zahl der Witwen nach 10 Jahren 600, nach 20 Jahren 1000, nach 30 Jahren 1200, nach 40 Jahren 1280, nach 45 Jahren 1294, nach 50 Jahren 1298 betragen, zu welchen später noch ein paar hinzukommen, und das wirkliche Maximum 1300 hervorbringen. Hier hätten wir demnach das Verhältniss der stehenden Ehen zu der höchsten Witwenzahl wie zwei zu eins.

In einer wirklichen Gesellschaft, wo die geforderten Bedingungen nicht in ihrer scharfen Strenge, sondern nur durchschnittlich gelten, wird es natürlich nicht so regelmässig hergehen können, wie in der fingirten. In jener werden jährlich nicht genau 130 neue Ehepaare beitreten, sondern in einem Jahre etwas mehr, in einem andern etwas weniger. Eben so werden die verschiedenen Altersstufen der neu beitretenden in einem Jahre etwas anders gemischt sein, als in einem andern, oder als in der idealen Gesellschaft angenommen war, zum Beispiel, das Durchschnittsalter der Männer, oder das der Frauen, oder der durchschnittliche Unterschied beider wird einmal etwas grösser, ein andermal etwas kleiner sein. Endlich wird auch von einer gegebenen Personenmenge das Absterben nicht \textit{genau} nach den Mortalitätstafeln erfolgen, sondern in einem Jahre werden diese etwas zu wenig, in einem andern etwas zu viel angeben. Der Erfolg von allem dem wird sein, dass in der wirklichen Gesellschaft zu einer Zeit die Zahl der Witwen etwas grösser sein wird, als in der idealen, zu einer andern etwas kleiner. Ein eigentlicher Beharrungszustand \textit{im strengsten Sinn} wird in jener niemals eintreten, sondern ein Fluctuiren um einen Mittel%
